\section{Planteamiento del problema}

En la prestación de servicios de soporte técnico y mantenimiento de equipos informáticos, la gestión documental de la recepción de dispositivos, el registro detallado del diagnóstico técnico y la comunicación continua sobre el avance de las órdenes de trabajo constituyen actividades operativas críticas para garantizar la transparencia y la calidad en la atención al cliente.

En el contexto de los prestadores de servicios de soporte técnico de la ciudad de Tarija, la ausencia de herramientas informáticas especializadas conduce a que las actividades de recepción y seguimiento se realicen habitualmente de manera empírica e informal. En la práctica cotidiana, los datos de los equipos recibidos y los diagnósticos iniciales suelen registrarse en boletas físicas de papel o en conversaciones dispersas a través de aplicaciones de mensajería instantánea. Esta dispersión documental dificulta la reconstrucción del historial de mantenimiento, complica la verificación de los componentes sustituidos y limita la capacidad de auditar las acciones realizadas entre el técnico y el usuario.

Asimismo, los clientes experimentan barreras para consultar de forma oportuna el estado actualizado de su reparación, dependiendo de llamadas telefónicas o visitas presenciales que interrumpen la labor técnica del personal en taller. Esta dinámica genera incertidumbre en los usuarios respecto a los tiempos de entrega y la fidelidad de los diagnósticos, planteando como hipótesis de trabajo la existencia de una desconfianza latente que requiere ser investigada empíricamente en el campo mediante entrevistas y encuestas estructuradas.